\makeatletter
\def\input@path{{./}{../}{../../}{../../../}{../../../../}{preamble/}{../preamble/}{../../preamble/}{../../../preamble/}}
\makeatother
\input{preamble_exams.tex}

\title{Grade 10 -- Term 1\\Comparing Decision Criteria: C8 Practice}
\author{}
\date{}

\begin{document}
\maketitle
\setcounter{tocdepth}{1}
\setcounter{secnumdepth}{2}
\phantomsection\label{toc}
\tableofcontents

\section*{Evaluated criteria}
\begin{enumerate}[label=\textbf{C\arabic*},start=8]
\item Compares the Maximax, Maximin, Minimax Regret, and EMV/Bayes decision criteria by contrasting their information requirements, risk orientations, recommended alternatives, and limitations.
\end{enumerate}

\section{Introduction}\label{sec:c8-introduction}
All problems use \textbf{profit payoffs}: larger values are better, and negative values represent losses. Alternatives $A_i$ occupy rows; mutually exclusive states of nature $S_j$ occupy columns. The payoff $P_{ay}(A_i,S_j)$ is the profit if $A_i$ is chosen and $S_j$ occurs. A decision is made before the state is known.

\subsection{Four ways to compare the same alternatives}
\textbf{Maximax} selects the alternative with the highest row maximum. It asks, ``Which alternative offers the greatest possible profit?'' \textbf{Maximin} selects the alternative with the highest row minimum. It asks, ``Which alternative has the best worst-case profit?'' Neither needs probabilities, even when probabilities are available.

\textbf{Minimax Regret} first compares alternatives within each state:
\[
b(S_j)=\max_i P_{ay}(A_i,S_j),\qquad
r(A_i,S_j)=b(S_j)-P_{ay}(A_i,S_j).
\]
Regret is the profit forgone compared with the best alternative for the state that actually occurs; it is not necessarily a cash loss. Select the alternative with the \emph{smallest maximum regret}. This criterion needs the full payoff table, but no probabilities.

\textbf{EMV/Bayes} selects the highest Expected Monetary Value:
\[
EMV(A_i)=\sum_{j=1}^{n}P(S_j)P_{ay}(A_i,S_j),\qquad
P(S_j)\geq0,\quad\sum_{j=1}^{n}P(S_j)=1.
\]
It requires a probability for every state as well as the payoffs. In these problems the state probabilities are the same for every alternative. Maximizing expected money expresses a risk-neutral orientation: it does not directly penalize uncertainty or protect against a large loss. An EMV is a weighted average, not a guaranteed profit or necessarily an attainable single outcome.

\begin{center}
\small
\begin{tabularx}{\textwidth}{@{}>{\raggedright\arraybackslash}p{0.15\textwidth}*{3}{>{\raggedright\arraybackslash}X}@{}}
\toprule
Criterion & Information and selection & Risk orientation & Main limitation\\
\midrule
Maximax & Payoffs; highest row maximum & Optimistic; pursues upside & Ignores other outcomes and their likelihoods.\\
Maximin & Payoffs; highest row minimum & Conservative; protects the worst payoff & Ignores upside and how unlikely the worst state may be.\\
Minimax Regret & Full payoff table; smallest row maximum regret & Cautious about missed opportunities & Ignores likelihoods; depends on the alternatives included.\\
EMV/Bayes & Payoffs and valid probabilities; highest EMV & Risk-neutral about money & Sensitive to probability estimates; may recommend an unaffordable downside.\\
\bottomrule
\end{tabularx}
\end{center}
Criteria can disagree because they compare different summaries of the \emph{same} information. Agreement does not mean that their reasoning is identical. Report all tied alternatives; do not invent a tie-breaking rule. Without probability information, EMV cannot be calculated unless an additional probability assumption is explicitly justified. Equal probabilities are not implied by uncertainty.
\secInicio[][sec:c8-introduction]

\section{Comparing Decision Criteria}\label{sec:c8-comparing}
\SectionProblemsTableOfContents{
\SectionProblemLink{sec:c8-worked}{Worked comparison: delivery expansion}\\
\SectionProblemLink{sec:c8-festival}{Festival catering: changing probabilities}\\
\SectionProblemLink{sec:c8-agreement}{Retail plans: agreement and ties}
}

\subsection{Worked Comparison: Delivery Expansion}\label{sec:c8-worked}
A delivery firm chooses its expansion plan. Quarterly profits are in thousand US dollars. Forecast probabilities are shown in the table.
\begin{center}
\begin{tabular}{lccc}\toprule
Alternative & Strong ($S_1$) & Stable ($S_2$) & Weak ($S_3$)\\
$P(S_j)$ & 0.20 & 0.60 & 0.20\\\midrule
Large fleet ($A_1$) & 100 & 20 & -20\\
Flexible fleet ($A_2$) & 75 & 50 & 10\\
Current fleet ($A_3$) & 40 & 40 & 35\\
Contract partners ($A_4$) & 60 & 65 & 20\\\bottomrule
\end{tabular}
\end{center}
\QuestionsTasks[sec:c8-comparing]{
{8}/{subsec:c8-worked-solution}/{Compare all four recommendations and explain why they differ.}
}
\subsubsection*{C8 -- Worked comparison}
\phantomsection\label{subsec:c8-worked-solution}
The statewise best payoffs are $(100,65,35)$. The regret table is
\begin{center}
\begin{tabular}{lrrrr}\toprule
Alternative & $S_1$ & $S_2$ & $S_3$ & Maximum regret\\\midrule
$A_1$ & 0 & 45 & 55 & 55\\
$A_2$ & 25 & 15 & 25 & 25\\
$A_3$ & 60 & 25 & 0 & 60\\
$A_4$ & 40 & 0 & 15 & 40\\\bottomrule
\end{tabular}
\end{center}
For example, $r(A_1,S_3)=35-(-20)=55$. The probabilities are valid since $0.20+0.60+0.20=1$. For the contract partners,
\[
EMV(A_4)=0.20(60)+0.60(65)+0.20(20)=55.
\]
\begin{center}
\begin{tabular}{lrrrr}\toprule
Alternative & Row maximum & Row minimum & Maximum regret & EMV\\\midrule
$A_1$ & 100 & -20 & 55 & 28\\
$A_2$ & 75 & 10 & 25 & 47\\
$A_3$ & 40 & 35 & 60 & 39\\
$A_4$ & 65 & 20 & 40 & 55\\\bottomrule
\end{tabular}
\end{center}
Maximax recommends $A_1$ for its upside of 100; Maximin recommends $A_3$ for its worst payoff of 35; Minimax Regret recommends $A_2$ for its maximum regret of only 25; EMV/Bayes recommends $A_4$ for its EMV of 55. The stable state has the greatest probability and $A_4$ performs best in that state. The other three criteria do not use those probabilities.

There is no calculation error in this disagreement. Each recommendation answers a different question. A firm unable to absorb a loss of 20 may reject $A_1$ despite its upside. $A_3$ protects profit but can miss 60 in the strong state. $A_2$ limits missed opportunity rather than guaranteeing the highest minimum profit. $A_4$ maximizes expected profit but can still earn only 20.
\secInicio[sec:c8-worked][sec:c8-comparing]

\subsection{Festival Catering: Changing Probabilities}\label{sec:c8-festival}
A caterer chooses a large or small contract before knowing attendance. Profits are in thousand US dollars.
\begin{center}
\begin{tabular}{lcc}\toprule
Alternative & High ($S_1$) & Low ($S_2$)\\\midrule
Large contract ($A_1$) & 80 & 10\\
Small contract ($A_2$) & 45 & 35\\\bottomrule
\end{tabular}
\end{center}
\begin{enumerate}
\item With $P(S_1)=0.30$ and $P(S_2)=0.70$, calculate row maxima, row minima, the regret table, maximum regrets, and EMVs. Summarize all four recommended alternatives in one table.
\item Repeat only the EMV calculations for $P(S_1)=0.60$ and $P(S_2)=0.40$. Explain why the other three recommendations stay unchanged.
\item Compare the reasoning of Maximin and Minimax Regret, even if their recommendations differ. Does limiting maximum regret necessarily protect the lowest profit?
\item Recommend a criterion for a caterer prioritizing expected profit and another for a caterer prioritizing the worst payoff. Cite the relevant numerical evidence and one limitation of each choice.
\end{enumerate}
\secInicio[sec:c8-festival][sec:c8-comparing]

\subsection{Retail Plans: Agreement and Ties}\label{sec:c8-agreement}
A retailer compares three plans. Profits are in thousand CNY.
\begin{center}
\begin{tabular}{lccc}\toprule
Alternative & Strong ($S_1$) & Normal ($S_2$) & Weak ($S_3$)\\
$P(S_j)$ & 0.25 & 0.50 & 0.25\\\midrule
Plan $A_1$ & 60 & 45 & 30\\
Plan $A_2$ & 50 & 40 & 20\\
Plan $A_3$ & 60 & 35 & 30\\\bottomrule
\end{tabular}
\end{center}
\begin{enumerate}
\item Apply all four criteria and report every tie. Use a regret table and a comparison table to support your conclusions.
\item Compare $A_1$ with each other alternative in every state. Explain why $A_1$ is never worse, yet some criteria still allow a tie.
\item Explain why agreement on $A_1$ would not make the four criteria interchangeable. Identify which summary each criterion uses.
\end{enumerate}
\secInicio[sec:c8-agreement][sec:c8-comparing]

\section{Contextual Practice}\label{sec:c8-contextual}
\SectionProblemsTableOfContents{
\SectionProblemLink{sec:c8-farm}{Urban farm crop mix}\\
\SectionProblemLink{sec:c8-investment}{Business reserve investment}\\
\SectionProblemLink{sec:c8-ferry}{Coastal ferry capacity}
}
For each problem, show row maxima and minima, statewise best payoffs, the regret table and maximum regrets, and all EMVs when probabilities are provided. Complete a comparison table with \textbf{criterion, recommended alternative, deciding value, and limitation}. Support explanations with values from the problem.

\subsection{Urban Farm Crop Mix}\label{sec:c8-farm}
An urban farm chooses a crop mix for the next season. Profits are in thousand US dollars; the estimates come from a market forecast.
\begin{center}
\begin{tabular}{lccc}\toprule
Alternative & Strong ($S_1$) & Stable ($S_2$) & Weak ($S_3$)\\
$P(S_j)$ & 0.50 & 0.30 & 0.20\\\midrule
Specialty crops ($A_1$) & 80 & 40 & -10\\
Balanced crops ($A_2$) & 60 & 55 & 20\\
Staple crops ($A_3$) & 40 & 45 & 35\\\bottomrule
\end{tabular}
\end{center}
\begin{enumerate}
\item Apply and compare all four criteria. Explain why the largest possible profit, the best worst-case profit, and the smallest maximum regret need not belong to one crop mix.
\item The farm must retain at least 30 in profit in every state to meet its obligations. Which alternatives meet that condition? Compare a recommendation under this constraint with the unrestricted EMV recommendation.
\item A revised forecast gives probabilities $(0.20,0.30,0.50)$. Recalculate EMVs. Explain what changes, what stays unchanged, and why reliable probability estimates matter.
\end{enumerate}
\secInicio[sec:c8-farm][sec:c8-contextual]

\subsection{Business Reserve Investment}\label{sec:c8-investment}
A business invests part of its reserve for one year. The table gives \emph{net gains}, in thousand US dollars, rather than final balances. All listed states are possible.
\begin{center}
\begin{tabular}{lccc}\toprule
Alternative & Growth ($S_1$) & Stable ($S_2$) & Recession ($S_3$)\\
$P(S_j)$ & 0.60 & 0.25 & 0.15\\\midrule
Growth fund ($A_1$) & 90 & 25 & -30\\
Balanced fund ($A_2$) & 60 & 40 & 5\\
Fixed deposit ($A_3$) & 25 & 25 & 25\\\bottomrule
\end{tabular}
\end{center}
\begin{enumerate}
\item Apply all four criteria and compare their recommendations. Explain any agreement using the deciding values, rather than claiming that the criteria have the same risk orientation.
\item Compare the growth fund's EMV with its possible recession loss. Explain why expected gain does not guarantee that the reserve is protected.
\item The treasurer prioritizes the highest expected gain; the owner requires a nonnegative gain in every state. Give each a justified recommendation, identifying the criterion or constraint used and its limitation.
\end{enumerate}
\secInicio[sec:c8-investment][sec:c8-contextual]

\subsection{Coastal Ferry Capacity}\label{sec:c8-ferry}
A ferry operator selects a seasonal capacity plan. Profits are in thousand US dollars. Initially, no reliable probabilities are available.
\begin{center}
\begin{tabular}{lccc}\toprule
Alternative & High ($S_1$) & Normal ($S_2$) & Low ($S_3$)\\\midrule
Add two ferries ($A_1$) & 100 & 35 & -25\\
Add one ferry ($A_2$) & 75 & 55 & 10\\
Lease peak capacity ($A_3$) & 55 & 48 & 30\\
Keep current fleet ($A_4$) & 30 & 30 & 30\\\bottomrule
\end{tabular}
\end{center}
\begin{enumerate}
\item Apply Maximax, Maximin, and Minimax Regret. Explain their disagreements and report any ties. State why EMV cannot yet be determined.
\item A later forecast supplies $P(S_1)=0.20$, $P(S_2)=0.30$, and $P(S_3)=0.50$. Check validity, calculate every EMV, and compare Bayes with the earlier recommendations.
\item The operator dislikes missing opportunities but has no trustworthy forecast. Which criterion fits that priority? Explain why it does not necessarily select the plan with the safest profit.
\end{enumerate}
\secInicio[sec:c8-ferry][sec:c8-contextual]

\section{Mixed Practice}\label{sec:c8-mixed}
\SectionProblemsTableOfContents{
\SectionProblemLink{sec:c8-match}{Choosing a criterion}\\
\SectionProblemLink{sec:c8-critique}{Evaluating decision claims}\\
\SectionProblemLink{sec:c8-report}{Decision committee report}
}

\subsection{Choosing a Criterion}\label{sec:c8-match}
For each situation, select the most appropriate of the four criteria \emph{for the stated priority}. Explain its information requirements, risk orientation, and one limitation. If another criterion is feasible, explain why it serves a different priority.
\begin{enumerate}
\item A start-up has a complete payoff table but no reliable probabilities. Its founder wants the largest possible profit and accepts the downside.
\item A family business has reliable probabilities but prioritizes the highest worst-case profit, even if expected profit is lower.
\item A delivery manager has no reliable probabilities and wants to minimize the greatest profit forgone relative to the best plan in each state.
\item A retailer makes many comparable seasonal decisions, trusts the state probabilities, and prioritizes average monetary profit over time.
\item A museum knows its payoffs but has no attendance probabilities. Its director requests the plan with the highest EMV. Identify the missing information and explain why choosing equal probabilities would introduce an assumption.
\end{enumerate}
\secInicio[sec:c8-match][sec:c8-mixed]

\subsection{Evaluating Decision Claims}\label{sec:c8-critique}
Decide whether each statement is valid. Correct invalid statements and support your response with a numerical example from this activity where possible.
\begin{enumerate}
\item ``Maximin and Minimax Regret must select the same alternative because both are cautious.''
\item ``If an alternative has an EMV of 55, it will earn 55 in the actual season.''
\item ``When probabilities change but payoffs stay fixed, all four recommendations must change.''
\item ``Regret of 25 means an actual monetary loss of 25.''
\item ``If Maximax and Bayes select the same alternative, Bayes is an optimistic criterion.''
\item ``A valid probability distribution is enough to make EMV the best criterion for every decision-maker.''
\end{enumerate}
\secInicio[sec:c8-critique][sec:c8-mixed]

\subsection{Decision Committee Report}\label{sec:c8-report}
Return to the delivery expansion problem in Section~\ref{sec:c8-worked}. Four committee members prioritize, respectively, the largest possible profit, the best worst-case profit, the smallest maximum missed opportunity, and the highest expected profit.
\begin{enumerate}
\item Match each member to a criterion and an alternative. Give the numerical deciding value and explain which information the member uses or ignores.
\item Write a short recommendation to the firm that compares at least two competing plans. State your chosen decision priority, use numerical evidence, and acknowledge one limitation. Explain why another member can reasonably disagree.
\item Suppose the forecast is withdrawn. Identify which calculations remain usable and which recommendation loses its information basis. Explain why this does not automatically make any one remaining criterion universally best.
\end{enumerate}
\textbf{Self-check.} A complete comparison identifies the available information, explains the decision-maker's priority, reports the recommended alternatives (including ties), and discusses the outcomes or uncertainty that each criterion leaves out.
\secInicio[sec:c8-report][sec:c8-mixed]
\end{document}

